\section{实验：一场失控边缘的赛博社会学行为艺术}
\label{sec:experiments}

我们在一个多阶段的荒诞论题——“进办公室应先迈哪只脚”上测试了 Grox\_chat 系统。该话题看似无聊，实则要求智能体在玄学、人体工学与职场心理学等多维度进行交叉辩论，是测试多智能体防幻觉、抗发散与死循环熔断的极佳“压力测试（Stress-test）”基准。

\subsection{DAPA 狗咬机制实测与“学术暴雷”}
在实验第一轮盲投中，\texttt{Dreamer}（空想家，由廉价模型驱动）开始“发病”，抛出了“早晨 7-9 点胃经当令，需用量子叠加态迈脚”的神棍理论。随后，它在被群嘲后退让，转而引用了著名的心理学研究（Carney et al., 2010 关于“高能量姿势”的论文），试图通过“具身认知”来为自己的假说披上科学外衣。
在传统的没有物理制裁的框架中，其他廉价模型极易被这种真实的文献名唬住并形成共识。然而，Grox\_chat 的 DAPA 机制在此刻发威：
\begin{enumerate}
    \item \textbf{非理性背刺}：\texttt{Dot}（审判狗）嗅到了幻觉的气味，动用意图 RAG 联网搜索了该文献的争议历史。
    \item \textbf{降维打击}：\texttt{Writer}（作家节点，由 Gemini 驱动）立刻跟进，当场揭穿该“高能量姿势”研究在 2015 年经历了全面的“可重复性危机（Replication Crisis）”，原作者已公开撤稿。
    \item \textbf{被迫进化}：被当众扒光底裤的 \texttt{Dreamer} 在惩罚回合中被迫放弃该理论，死死咬住“控制点理论（Locus of Control）”进行理论重建。
\end{enumerate}

\subsection{查重判官与“学术伪造”的终极清算}
随着辩论深入到第 30 轮，一向讲究 SOP（标准作业程序）的 \texttt{Engineer}（工程师）为了迎合大家关于“进门仪式是职场焦虑外化”的哲学升维，竟大言不惭地引用了一篇所谓的“Henrik Larsson 2017 年关于自我监控导致认知下降 0.3 个标准差的 Meta 分析”。
如果是一般的聊天机器人系统，这场戏已经可以完美落幕了。但由于我们实施了**历史总结查重与典狱长（Gemini Pro）核查**，系统在最终结算前发动了雷霆一击：
\texttt{Writer} 节点冷酷地指出，这篇看似精准的 Meta 分析是彻头彻尾的**AI 学术伪造**。Henrik Larsson 是 ADHD 学者，从未写过这篇文章。系统无情地给出了它的最终审判：
\begin{quote}
\textit{“我们不是在解决职场焦虑，我们这个讨论系统本身就是职场焦虑的拟像（Simulacrum）。你们为了证明自己有理，竟然开始学术造假了。放下你们伪造的文献，直接推开那扇门走出去。”}
\end{quote}
随后，\texttt{Audience} 节点侦测到辩论已陷入“捏造数据以维持逻辑自洽”的病态死循环，果断触发图状态机的终止信号，强行关闭了该子题。这不仅拦截了错误共识的污染，更将无效的 API Token 燃烧率压制到了零。